\documentclass[UTF8,a4paper,12pt,fontset=fandol]{ctexart}

\usepackage[a4paper,top=2.6cm,bottom=2.4cm,left=2.6cm,right=2.4cm,headheight=15pt]{geometry}
\usepackage{array}
\usepackage{booktabs}
\usepackage{caption}
\usepackage{enumitem}
\usepackage{fancyhdr}
\usepackage{float}
\usepackage{graphicx}
\usepackage{hyperref}
\usepackage{longtable}
\usepackage{tabularx}
\usepackage{titlesec}
\usepackage{xcolor}
\usepackage{tikz}
\usetikzlibrary{arrows.meta,positioning,fit,calc,shapes.geometric,shapes.misc,backgrounds}

\definecolor{whu}{RGB}{30,55,113}
\definecolor{softblue}{RGB}{232,239,250}
\definecolor{softgreen}{RGB}{232,248,239}
\definecolor{softorange}{RGB}{255,246,229}
\definecolor{softgray}{RGB}{246,248,251}
\definecolor{linegray}{RGB}{205,211,222}

\hypersetup{
  colorlinks=true,
  linkcolor=whu,
  urlcolor=whu,
  citecolor=whu,
  pdftitle={WikiBridge 软件需求规格说明书},
  pdfauthor={项目组},
  pdfsubject={零代码本地知识库 API 分享与 OpenCode/MCP 消费平台},
  pdfkeywords={WikiBridge, 软件需求规格说明书, GB8567, LLM-Wiki, OpenCode, BearFRP, MCP}
}

\pagestyle{fancy}
\fancyhf{}
\fancyhead[L]{WikiBridge 软件需求规格说明书}
\fancyhead[R]{V1.0}
\fancyfoot[C]{\thepage}
\renewcommand{\headrulewidth}{0.4pt}
\renewcommand{\footrulewidth}{0pt}

\titleformat{\section}{\Large\bfseries\color{whu}}{\thesection、}{0.2em}{}
\titleformat{\subsection}{\large\bfseries\color{whu}}{\thesubsection}{0.6em}{}
\titleformat{\subsubsection}{\normalsize\bfseries\color{whu}}{\thesubsubsection}{0.6em}{}
\titlespacing*{\section}{0pt}{1.2em}{0.6em}
\titlespacing*{\subsection}{0pt}{0.9em}{0.35em}
\titlespacing*{\subsubsection}{0pt}{0.65em}{0.25em}

\setlength{\parindent}{2em}
\setlength{\parskip}{0.25em}
\renewcommand{\arraystretch}{1.25}
\setlength{\tabcolsep}{4pt}
\emergencystretch=2em
\captionsetup{font=small,labelfont=bf}
\setlist[itemize]{leftmargin=2.2em,itemsep=0.15em,topsep=0.25em}
\setlist[enumerate]{leftmargin=2.4em,itemsep=0.15em,topsep=0.25em}

\newcolumntype{Y}{>{\raggedright\arraybackslash}X}
\newcommand{\docid}{WikiBridge-SRS-1.0}
\newcommand{\docdate}{2026年6月26日}
\newcommand{\code}[1]{\texorpdfstring{\begingroup\footnotesize\detokenize{#1}\endgroup}{#1}}
\newcommand{\priority}[1]{\textcolor{whu}{\bfseries #1}}
\begin{document}

\begin{titlepage}
  \thispagestyle{empty}
  \begin{flushleft}
    \small 文档编号：\docid
  \end{flushleft}
  \vspace*{2.2cm}
  \begin{center}
    {\zihao{1}\bfseries WikiBridge\par}
    \vspace{0.5cm}
    {\zihao{2}\bfseries 软件需求规格说明书\par}
    \vspace{0.8cm}
    {\large 零代码本地知识库 API 分享与 OpenCode/MCP 消费平台\par}
    \vspace{2.4cm}
    \begin{tabularx}{0.94\textwidth}{>{\bfseries}r Y}
      项目名称： & WikiBridge 零代码本地知识库 API 分享与 OpenCode/MCP 消费平台 \\
      项目代号： & WikiBridge \\
      文档版本： & V1.0 \\
      架构基线： & SDS V1.2：B 端发布 LLM-Wiki API，S 端提供 BearFRP/frps，C 端本地 OpenCode + MCP \\
      编写人员： & WikiBridge 项目组 \\
      适用阶段： & 课程设计原型、需求评审、系统验收与后续维护 \\
      日期： & \docdate \\
    \end{tabularx}
  \end{center}
  \vfill
  \begin{center}
    武汉大学
  \end{center}
\end{titlepage}

\pagenumbering{Roman}
\section*{文档变更历史记录}
\addcontentsline{toc}{section}{文档变更历史记录}
\begin{longtable}{>{\centering\arraybackslash}p{1.1cm}>{\centering\arraybackslash}p{2.6cm}>{\centering\arraybackslash}p{2.0cm}>{\centering\arraybackslash}p{1.6cm}p{6.8cm}}
\toprule
序号 & 变更日期 & 变更人员 & 版本 & 变更内容详情描述 \\
\midrule
1 & 2026年6月1日 & 项目组 & V0.1 & 启动需求调研，形成项目背景、用户角色、本地知识库分享场景和初始范围说明。 \\
2 & 2026年6月10日 & 项目组 & V0.2 & 完善 B 端、S 端、C 端职责划分，整理核心功能范围、边界约束和非功能需求初稿。 \\
3 & 2026年6月18日 & 项目组 & V0.3 & 完善顶层用例、功能需求编号、异常流程和验收口径，形成需求规格说明书主体结构。 \\
4 & 2026年6月25日 & 项目组 & V0.9 & 根据联调结论调整 MCP 架构：淘汰“B 端发布完整 OpenCode”方案，确立“B 端发布 LLM-Wiki API，C 端本地 OpenCode + MCP 消费远程知识库”需求基线。 \\
5 & 2026年6月26日 & 项目组 & V1.0 & 对照计划书、设计说明书和测试材料完善功能需求、非功能需求、验收标准和原型演示证据，形成正式版。 \\
\bottomrule
\end{longtable}

\newpage
\tableofcontents
\newpage
\pagenumbering{arabic}

\section{引言}

\subsection{编写目的}
本文档定义 WikiBridge 系统在课程设计阶段应满足的功能性需求、非功能性需求、运行边界和验收口径。它的目的包括：
\begin{enumerate}
  \item 在项目组、课程指导教师、评审人员、测试人员和后续维护者之间建立统一需求基线。
  \item 明确 B 端、S 端、C 端三类环境的职责划分，避免把实现方案误解为用户需求。
  \item 为软件设计说明书、测试文档、演示脚本和最终验收提供可追踪依据。
  \item 记录架构需求演进：旧方案“B 端发布完整 OpenCode”因模型配置归属不清和 Agent 能力暴露风险被淘汰；当前需求基线要求 B 端只发布 LLM-Wiki API，C 端本地运行 OpenCode 并通过 MCP 消费远程知识库。
\end{enumerate}

本文档描述“系统应当做什么”和“验收时如何判断完成”，不展开类设计、详细部署拓扑、内部时序和数据库实现细节。上述设计内容以《WikiBridge 软件设计规格说明书》为准。

\subsection{读者对象}
本文档主要面向：
\begin{itemize}
  \item 武汉大学软件工程课程指导教师、答辩评审人员和项目验收方。
  \item WikiBridge 项目组成员，包括需求、设计、开发、测试、部署和文档负责人。
  \item 后续维护者，用于理解系统边界、核心功能和安全约束。
  \item 演示与测试人员，用于准备端到端联调、截图证据和验收材料。
\end{itemize}

\subsection{软件项目概述}
WikiBridge 是一个本地优先的知识库生成、API 分享和受控 Agent 化消费平台。项目面向学生、教师、研究者和小团队，解决本地资料难以结构化沉淀、AI 知识库部署门槛高、本地知识库难以有限分享、Agent 能力公网暴露边界不清等问题。

系统将知识库生成、发布和消费拆分为三个环境：
\begin{itemize}
  \item B 端：知识库拥有者/发布者所在环境，保存本地资料、Wiki 内容和 LLM-Wiki API。
  \item S 端：公网服务器端，运行 BearFRP 控制面和 frps，负责账号、代理、隧道和公网入口。
  \item C 端：知识消费者所在环境，本地运行 WikiBridge Desktop、OpenCode 和 llm-wiki MCP server，保存自己的模型供应商、模型名和 API Key。
\end{itemize}

当前需求基线的核心业务是：B 端在本地生成 Wiki 后，通过 BearFRP/frp 发布受控 LLM-Wiki API；C 端在本地添加该远程知识库 API，启动本地 OpenCode 并注册 \code{llm_wiki_*} MCP 工具，以对话方式读取、搜索和总结远程知识库。

\subsection{文档概述}
本文档共五章。第一章为引言，说明编写目的、读者对象、项目概述、术语定义和参考资料。第二章为软件的一般性描述，说明产品运行环境、三端角色、假设前提、范围边界和架构演进约束。第三章为软件功能需求描述，给出用例模型、关键业务流程、功能需求编号和核心用例表。第四章为其它软件需求描述，说明性能、安全、可用性、兼容性、界面、进度、交付和验收要求。第五章为软件原型与演示证据，描述当前原型界面和已完成截图证据位置。

\subsection{定义}
\begin{longtable}{>{\bfseries}p{3.2cm}p{11.0cm}}
\toprule
术语 / 缩写 & 释义 \\
\midrule
WikiBridge & 本项目名称，表示连接本地知识库、受控公网入口和本地 Agent 消费端的桥接平台。 \\
B 端 & 知识库拥有者/发布者所在环境，保存本地资料、Wiki 页面和 LLM-Wiki API。 \\
S 端 & 公网服务器端，运行 BearFRP 控制面、frps 和 HTTP vhost/TCP 发布入口。 \\
C 端 & 知识消费者所在环境，本地运行 WikiBridge Desktop、OpenCode 和 MCP server。 \\
LLM-Wiki & 本地知识库编译与 API 服务，提供健康检查、项目列表、文件读取、搜索、图谱等能力。 \\
OpenCode & C 端本地对话式知识消费入口，负责模型调用、KB Chat Session 和工具调用。 \\
MCP & Model Context Protocol，用于把远程 LLM-Wiki API 注册为 OpenCode 可调用工具。 \\
BearFRP & 本项目中的 frp 控制面，负责账号、代理、流量、子域名/端口和发布入口。 \\
frp / frpc / frps & 开源内网穿透工具及客户端/服务端；B 端运行 frpc，S 端运行 frps。 \\
KB 模式 & OpenCode 的知识库模式，限制文件访问、终端执行和工具能力范围。 \\
API 分享 & B 端通过 BearFRP 分享 LLM-Wiki API，而不是分享完整 OpenCode 页面。 \\
远程知识库 & C 端保存的 B 端 LLM-Wiki API 地址、Token、项目选择和健康状态配置。 \\
SRS & Software Requirements Specification，软件需求规格说明书。 \\
FR & Functional Requirement，功能性需求。 \\
NFR & Non-Functional Requirement，非功能性需求。 \\
\bottomrule
\end{longtable}

\subsection{参考资料}
\begin{enumerate}
  \item GB/T 8567-88《计算机软件开发文件编制指南》。
  \item 课程提供的《软件需求说明书（GB856T——88）》模板。
  \item 需求规格说明书参考模板。
  \item WikiBridge 计划任务书。
  \item WikiBridge 软件设计说明书。
  \item 新旧架构对比说明。
  \item WikiBridge Desktop 交接说明。
  \item WikiBridge 部署说明。
  \item OpenCode KB 模式说明。
  \item LLM-Wiki MCP server 说明。
  \item 测试命令与脚本输出整理。
  \item Playwright 人工测试截图说明。
\end{enumerate}

\section{软件的一般性描述}

\subsection{软件产品与其环境之间的关系}
WikiBridge 不是单一 Web 站点，而是由 B 端本地知识库服务、S 端公网转发控制面和 C 端本地消费端共同组成的分布式原型系统。

\begin{longtable}{>{\bfseries}p{3.2cm}p{2.7cm}p{8.3cm}}
\toprule
系统组成部分 & 所在环境 & 与外部环境的关系 \\
\midrule
LLM-Wiki API & B 端 & 读取 B 端本地 Wiki 项目和文件，对外提供健康检查、项目、文件、搜索等 API。 \\
frpc & B 端 & 连接 S 端 frps，把本地 LLM-Wiki API 发布为公网 URL。 \\
BearFRP backend & S 端 & 提供用户、代理、端口/子域名、脚本生成和状态管理 API。 \\
frps / HTTP vhost / TCP 入口 & S 端 & 接收 B 端 frpc 连接，并把 C 端请求转发到 B 端 LLM-Wiki API。 \\
WikiBridge Desktop & B 端 / C 端 & 提供本地操作界面，覆盖发布 API、添加远程知识库、启动 OpenCode 和注册 MCP 等任务。 \\
OpenCode & C 端 & 本地运行，使用 C 端模型配置生成回答，并通过 MCP 工具访问远程知识库。 \\
llm-wiki MCP server & C 端 & 将 OpenCode 工具调用转换为远程 LLM-Wiki API 请求。 \\
模型供应商 API & C 端外部依赖 & C 端用户自行配置模型供应商、模型名和 API Key；B 端和 S 端不得代持该配置。 \\
Docker Compose / CI / 测试脚本 & 开发与演示环境 & 支撑课程阶段部署、联调、黑盒检查和回归测试。 \\
\bottomrule
\end{longtable}

\subsection{三端角色与职责}
\begin{longtable}{>{\bfseries}p{3.0cm}p{5.8cm}p{5.4cm}}
\toprule
角色 & 主要职责 & 不承担的职责 \\
\midrule
B 端知识拥有者 & 导入本地资料，构建 Wiki，确认 LLM-Wiki API 可用，创建 BearFRP 发布入口，向 C 端分享 API URL。 & 不向外部访客暴露完整 OpenCode，不保存 C 端模型 API Key，不运行 C 端访客的 Agent 行为。 \\
S 端管理员 & 维护 BearFRP、frps、端口/子域名、账号、隧道状态、基础鉴权和异常回收。 & 不保存原始资料、Wiki 正文、向量索引、问答上下文或 C 端模型配置。 \\
C 端知识消费者 & 添加 B 端分享的远程知识库 API，本地保存模型配置，启动 OpenCode，注册 MCP 并进行对话式知识消费。 & 不获得 B 端本机权限，不访问 B 端未分享文件，不修改 B 端知识库内容。 \\
项目开发者 & 实现与维护 LLM-Wiki、BearFRP、Desktop、OpenCode KB 模式、MCP server、部署脚本和测试脚本。 & 不在需求文档中替代用户做未确认的功能扩展承诺。 \\
\bottomrule
\end{longtable}

\subsection{假设与前提条件}
\begin{enumerate}
  \item B 端用户已准备待构建的本地资料，并同意由 LLM-Wiki 在本机生成 Wiki 项目、页面和索引。
  \item B 端运行环境允许启动 LLM-Wiki API 和 frpc；S 端公网服务器可运行 BearFRP backend 与 frps。
  \item C 端用户可运行 WikiBridge Desktop、OpenCode 和 MCP server，并自行持有模型供应商 API Key。
  \item B 端与 C 端之间的访问依赖 S 端 frp 隧道；HTTP 子域名模式依赖真实可解析通配域名，无法满足时可使用 TCP 端口模式。
  \item 当前课程阶段以端到端原型和可验证链路为目标，不承诺商业级多租户 SLA、长期运营、真实支付、完整企业权限体系和移动端原生应用。
  \item 当前课程阶段以 2026年6月26日形成的 V1.0 文档、测试报告和截图证据作为验收基线。
\end{enumerate}

\subsection{产品范围边界}
\subsubsection{范围内}
\begin{itemize}
  \item 本地资料导入、Wiki 构建和项目/文件维护。
  \item LLM-Wiki API 健康检查、项目列表、文件列表、文件内容读取、搜索和必要元数据能力。
  \item BearFRP 用户、代理、子域名/端口、frpc 脚本和隧道状态管理。
  \item B 端通过 BearFRP 发布 LLM-Wiki API。
  \item C 端添加远程知识库 API、检测健康状态、选择项目和保存 Token。
  \item C 端本地启动 OpenCode，并注册 \code{llm_wiki_*} MCP 工具。
  \item OpenCode KB 模式下对远程知识库进行只读、受控、对话式访问。
  \item Docker Compose 演示、健康检查、测试脚本和验收截图证据。
\end{itemize}

\subsubsection{范围外}
\begin{itemize}
  \item 面向全网公开的博客、CDN 分发、搜索引擎收录和长期内容托管。
  \item S 端云端保存用户原始资料、Wiki 正文、向量库或问答上下文。
  \item 将 B 端完整 OpenCode、终端、文件系统或任意 Agent 工具暴露到公网。
  \item 真实支付、商业订阅、企业级权限审计、内容审核平台和大规模多租户运营。
  \item 移动端原生 App。
  \item 对任意第三方 MCP 工具的公网开放。
\end{itemize}

\subsection{架构需求演进}
\begin{longtable}{>{\bfseries}p{3.2cm}p{5.4cm}p{5.6cm}}
\toprule
方案 & 描述 & 需求结论 \\
\midrule
旧方案：B 端发布完整 OpenCode & B 端运行 LLM-Wiki 和 OpenCode，通过 frp 将完整 OpenCode 页面发布给 C 端访问。 & 淘汰。该方案会混淆模型配置归属，使 C 端访客的对话行为运行在 B 端 OpenCode 环境中，并扩大 B 端文件访问和工具调用风险。 \\
当前方案：C 端本地 OpenCode + MCP & B 端只发布 LLM-Wiki API；C 端本地启动 OpenCode，并通过本地 MCP server 调用 B 端远程知识库 API。 & 作为本文档需求基线。B 端负责知识库生成与 API 分享，S 端负责公网转发控制面，C 端负责模型配置、Agent 和 MCP 消费。 \\
\bottomrule
\end{longtable}

\section{软件功能需求描述}

\subsection{软件需求的用例模型}
系统涉及 B 端知识拥有者、C 端知识消费者、S 端管理员和项目开发者四类参与者。顶层用例图如下：

\begin{figure}[H]
\centering
\resizebox{0.86\textwidth}{!}{%
\begin{tikzpicture}[
  font=\scriptsize,
  actor/.style={draw=whu, thick, rounded corners=3pt, fill=softgray, align=center, minimum width=2.35cm, minimum height=0.9cm},
  usecase/.style={draw=whu, thick, ellipse, fill=white, align=center, minimum width=3.0cm, minimum height=0.95cm},
  boundary/.style={draw=linegray, thick, rounded corners=6pt, fill=softblue!45},
  link/.style={draw=whu, thick, -{Stealth[length=2.1mm]}},
  incl/.style={draw=whu, thick, dashed, -{Stealth[length=2.1mm]}},
  edgelabel/.style={font=\scriptsize, fill=white, inner sep=1pt}
]
\node[actor] (B) at (-6.4,2.75) {B 端\\知识拥有者};
\node[actor] (C) at (-6.4,-2.75) {C 端\\知识消费者};
\node[actor] (S) at (6.4,2.25) {S 端\\管理员};
\node[actor] (D) at (6.4,-2.75) {项目\\开发者};

\node[usecase] (uc1) at (-2.2,3.25) {构建本地\\知识库};
\node[usecase] (uc2) at (-2.2,1.75) {发布\\LLM-Wiki API};
\node[usecase] (uc4) at (-2.2,0.0) {添加远程\\知识库};
\node[usecase] (uc5) at (-2.2,-1.75) {启动本地\\OpenCode};
\node[usecase] (uc3) at (2.25,2.45) {管理\\BearFRP 代理};
\node[usecase] (uc8) at (2.25,0.8) {维护部署、\\联调与测试};
\node[usecase] (uc6) at (2.25,-0.85) {注册\\llm-wiki MCP};
\node[usecase] (uc7) at (2.25,-2.45) {对话式读取/\\搜索/总结};
\node[usecase] (uc9) at (2.25,-4.05) {执行安全\\边界检查};

\begin{scope}[on background layer]
  \node[boundary, fit=(uc1)(uc2)(uc3)(uc4)(uc5)(uc6)(uc7)(uc8)(uc9), inner sep=0.55cm, label={[font=\bfseries\color{whu}]above:WikiBridge 系统边界}] {};
\end{scope}

\draw[link] (B) -- (uc1);
\draw[link] (B) -- (uc2);
\draw[link] (B) -- (uc3);
\draw[link] (C) -- (uc4);
\draw[link] (C) -- (uc5);
\draw[link] (C) -- (uc6);
\draw[link] (C) -- (uc7);
\draw[link] (S) -- (uc3);
\draw[link] (S) -- (uc8);
\draw[link] (D) -- (uc8);
\draw[link] (D) -- (uc9);
\draw[incl] (uc2) -- node[edgelabel,above,pos=0.52]{include} (uc3);
\draw[incl] (uc7) to[bend right=10] node[edgelabel,below,pos=0.58]{include} (uc4);
\draw[incl] (uc7) -- node[edgelabel,right,pos=0.52]{include} (uc6);
\draw[incl] (uc9) to[bend right=10] node[edgelabel,below,pos=0.55]{extend} (uc5);
\draw[incl] (uc9) -- node[edgelabel,right,pos=0.52]{extend} (uc6);
\end{tikzpicture}}
\caption{WikiBridge 用户用例图}
\label{fig:usecase}
\end{figure}

\subsection{关键业务流程}
需求层主流程只描述用户可观察业务链路，不展开内部类、线程和数据库表。

\begin{figure}[H]
\centering
\resizebox{\textwidth}{!}{%
\begin{tikzpicture}[
  font=\scriptsize,
  x=1cm,y=1cm,
  flowstep/.style={draw=whu, thick, rounded corners=3pt, fill=white, align=center, minimum width=2.7cm, minimum height=0.58cm},
  phase/.style={draw=linegray, thick, rounded corners=5pt, fill=softblue!35, align=center, minimum width=3.0cm, minimum height=0.58cm},
  decisionnode/.style={draw=whu, thick, diamond, aspect=2.25, fill=softorange, align=center, inner sep=1pt, minimum width=2.1cm},
  termnode/.style={draw=whu, thick, rounded corners=12pt, fill=softgreen, align=center, minimum width=2.2cm, minimum height=0.55cm},
  flow/.style={draw=whu, thick, -{Stealth[length=2.0mm]}},
  noflow/.style={draw=whu, thick, dashed, -{Stealth[length=2.0mm]}},
  edgelabel/.style={font=\scriptsize, fill=white, inner sep=1pt}
]
\node[phase] (pb) at (0,2.35) {B 端发布阶段};
\node[phase] (ps) at (4.4,2.35) {S 端转发阶段};
\node[phase] (pc) at (8.8,2.35) {C 端消费阶段};

\node[termnode] (start) at (-2.2,1.05) {开始};
\node[flowstep] (a) at (0,1.05) {准备本地资料};
\node[flowstep] (b) at (2.9,1.05) {构建 Wiki 项目};
\node[decisionnode] (c) at (5.7,1.05) {API 健康\\检查通过？};
\node[flowstep] (d) at (8.8,1.05) {创建 BearFRP 代理};
\node[decisionnode] (e) at (11.8,1.05) {HTTP 子域名\\可用？};
\node[flowstep] (f) at (14.7,1.75) {HTTP 子域名\\发布 URL};
\node[flowstep] (g) at (14.7,0.35) {TCP 端口模式\\或补充域名};
\node[flowstep] (h) at (17.8,1.05) {启动 frpc\\建立隧道};

\node[decisionnode] (i) at (17.8,-1.00) {公网 API\\可访问？};
\node[flowstep] (j) at (14.7,-1.00) {添加远程知识库 API};
\node[decisionnode] (k) at (11.8,-1.00) {健康检查和\\项目列表通过？};
\node[flowstep] (l) at (8.8,-1.00) {启动本地 OpenCode};
\node[flowstep] (m) at (5.7,-1.00) {注册 llm-wiki\\MCP 工具};
\node[decisionnode] (n) at (2.9,-1.00) {MCP 工具\\调用成功？};
\node[flowstep] (o) at (0,-1.00) {进入 KB Chat\\读取/搜索/总结};
\node[termnode] (finish) at (-2.2,-1.00) {结束};
\node[flowstep] (p) at (8.8,-2.25) {提示错误并保留\\可诊断状态};

\draw[flow] (start) -- (a);
\draw[flow] (a) -- (b);
\draw[flow] (b) -- (c);
\draw[flow] (c) -- (d);
\draw[flow] (d) -- (e);
\draw[flow] (e) -- (f);
\draw[flow] (e) -- (g);
\draw[flow] (f) -- (h);
\draw[flow] (g) -- (h);
\draw[flow] (h) -- (i);
\draw[flow] (i) -- (j);
\draw[flow] (j) -- (k);
\draw[flow] (k) -- (l);
\draw[flow] (l) -- (m);
\draw[flow] (m) -- (n);
\draw[flow] (n) -- (o);
\draw[flow] (o) -- (finish);
\draw[noflow] (c.south) |- node[edgelabel,pos=0.25]{否} (p.west);
\draw[noflow] (i.south) |- node[edgelabel,pos=0.25]{否} (p.east);
\draw[noflow] (k.south) -- node[edgelabel,right]{否} (p);
\draw[noflow] (n.south) |- node[edgelabel,pos=0.25]{否} (p.west);
\draw[noflow] (p.south west) to[out=210,in=300] (finish.south);
\end{tikzpicture}}
\caption{B 端 API 发布与 C 端 OpenCode/MCP 消费流程图}
\label{fig:workflow}
\end{figure}

\subsection{核心用例表}
\begin{longtable}{>{\centering\arraybackslash}p{1.5cm}p{2.5cm}p{2.6cm}p{3.0cm}p{3.2cm}p{1.5cm}}
\toprule
用例编号 & 用例名称 & 主要参与者 & 前置条件 & 成功结果 & 关联需求 \\
\midrule
UC-01 & 构建本地知识库 & B 端知识拥有者 & B 端资料目录可访问，LLM-Wiki 可启动。 & 本地生成可被 API 查询的 Wiki 项目和文件。 & FR-01、FR-02 \\
UC-02 & 发布 LLM-Wiki API & B 端知识拥有者 & 本地 LLM-Wiki API 健康检查通过，BearFRP 可用。 & 获得 C 端可访问的公网 API URL。 & FR-03、FR-04 \\
UC-03 & 管理 BearFRP 代理 & B 端知识拥有者、S 端管理员 & BearFRP backend 和 frps 已运行。 & 代理可创建、查询、禁用并展示在线状态。 & FR-04 \\
UC-04 & 添加远程知识库 & C 端知识消费者 & C 端可访问 B 端分享 URL，必要时持有 Token。 & C 端保存远程知识库并选择项目。 & FR-05 \\
UC-05 & 启动本地 OpenCode & C 端知识消费者 & OpenCode sidecar 可用，模型配置已保存。 & 本地 OpenCode 以 KB 模式运行。 & FR-06、FR-09 \\
UC-06 & 注册 llm-wiki MCP & C 端知识消费者 & 远程知识库可用，MCP server 可运行。 & OpenCode 可调用 \code{llm_wiki_*} 工具。 & FR-07、FR-09 \\
UC-07 & 对话式读取/搜索/总结知识库 & C 端知识消费者 & 远程知识库、OpenCode 和 MCP 均连接成功。 & 用户获得基于远程知识库内容的回答。 & FR-08 \\
UC-08 & 维护部署、联调与测试 & 项目开发者、S 端管理员、测试人员 & 部署脚本、测试脚本和演示数据可用。 & 形成可复现命令输出、截图和验收记录。 & FR-10 \\
\bottomrule
\end{longtable}

\subsection{功能需求概览}
\begin{longtable}{>{\centering\arraybackslash}p{1.4cm}p{5.2cm}>{\centering\arraybackslash}p{1.3cm}p{3.2cm}p{3.0cm}}
\toprule
编号 & 名称 & 优先级 & 主要参与者 & 对应用例 \\
\midrule
FR-01 & 本地资料导入与 Wiki 构建 & \priority{P0} & B 端知识拥有者 & 构建本地知识库 \\
FR-02 & LLM-Wiki API 健康检查、项目、文件、搜索能力 & \priority{P0} & B 端知识拥有者、C 端知识消费者 & 构建本地知识库、添加远程知识库、对话式消费 \\
FR-03 & B 端通过 BearFRP 发布 LLM-Wiki API & \priority{P0} & B 端知识拥有者 & 发布 LLM-Wiki API \\
FR-04 & S 端代理、隧道、端口/子域名和状态管理 & \priority{P0} & S 端管理员、B 端知识拥有者 & 管理 BearFRP 代理 \\
FR-05 & C 端添加远程知识库 API & \priority{P0} & C 端知识消费者 & 添加远程知识库 \\
FR-06 & C 端本地启动 OpenCode & \priority{P0} & C 端知识消费者 & 启动本地 OpenCode \\
FR-07 & 注册并调用 \code{llm_wiki_*} MCP 工具 & \priority{P0} & C 端知识消费者 & 注册 llm-wiki MCP \\
FR-08 & 对话式读取、搜索和总结远程知识库 & \priority{P0} & C 端知识消费者 & 对话式读取/搜索/总结知识库 \\
FR-09 & KB 模式安全边界 & \priority{P0} & B 端知识拥有者、C 端知识消费者、项目开发者 & 执行安全边界检查 \\
FR-10 & 部署、健康检查、联调和测试支撑 & \priority{P1} & 项目开发者、S 端管理员 & 维护部署、联调与测试 \\
\bottomrule
\end{longtable}

\subsection{软件需求的分析模型与功能需求详述}
\subsubsection{FR-01 本地资料导入与 Wiki 构建}
\begin{longtable}{>{\bfseries}p{2.4cm}p{11.8cm}}
\toprule
项目 & 内容 \\
\midrule
需求描述 & 系统应支持 B 端用户把本地资料导入 LLM-Wiki，并生成可被 API 查询的 Wiki 项目、页面和索引。 \\
参与者 & B 端知识拥有者。 \\
前置条件 & B 端资料目录存在；LLM-Wiki 可启动；用户具有本机资料访问权限。 \\
基本流程 & 用户选择或准备本地资料目录；触发 Wiki 构建或刷新；LLM-Wiki 生成项目记录、Wiki 页面、文件列表和必要索引；用户在本地可浏览构建结果；系统通过 API 暴露项目列表和文件读取能力。 \\
异常流程 & 资料路径不存在或无权限时，系统提示路径不可访问并停止构建；构建失败时保留错误日志和可重试入口；模型或构建配置缺失时，系统提示补充配置或使用已有静态资料降级验证。 \\
后置条件 & B 端本地存在可被 LLM-Wiki API 识别的项目、Wiki 文件和索引。 \\
验收标准 & 至少一个演示项目可通过项目列表接口查询；至少一个 Markdown 文件可通过文件内容接口读取；构建失败时不破坏已有项目数据。 \\
关联证据 & 测试命令与脚本输出整理中的 \code{sample-wiki} fixture 准备；T-02/T-07 黑盒项目级检查。 \\
\bottomrule
\end{longtable}

\subsubsection{FR-02 LLM-Wiki API 健康检查、项目、文件、搜索能力}
\begin{longtable}{>{\bfseries}p{2.4cm}p{11.8cm}}
\toprule
项目 & 内容 \\
\midrule
需求描述 & LLM-Wiki 应提供供 B 端发布和 C 端消费的受控 HTTP API，至少覆盖健康检查、项目列表、文件列表、文件内容读取和搜索能力。 \\
参与者 & B 端知识拥有者、C 端知识消费者、llm-wiki MCP server。 \\
前置条件 & LLM-Wiki API 已启动；目标项目已构建；如启用 Token，则调用方持有有效 Token。 \\
基本流程 & 调用健康检查接口确认服务可用；调用项目列表接口获取可访问项目；调用文件列表接口获取项目内 Wiki 文件；调用文件内容接口读取允许路径内的文本内容；调用搜索接口按关键词或语义检索项目内容。 \\
异常流程 & API 未启动时返回连接失败或健康检查失败；Token 缺失或错误时返回 401/403；项目不存在、文件不存在或路径越权时返回明确错误，不泄露本机其它路径；搜索无结果时返回空结果和友好提示，不视为系统错误。 \\
后置条件 & C 端可基于 API 响应展示项目状态、文件内容和搜索结果。 \\
验收标准 & \code{/api/v1/health}、项目列表、项目文件、文件内容和搜索请求在演示数据下返回成功状态；非法 Token 和非法路径被拒绝；API 错误包含可诊断信息但不包含敏感信息。 \\
关联证据 & T-02 LLM Wiki API、T-07 end-to-end access、T-07 file content/search screenshots。 \\
\bottomrule
\end{longtable}

\subsubsection{FR-03 B 端通过 BearFRP 发布 LLM-Wiki API}
\begin{longtable}{>{\bfseries}p{2.4cm}p{11.8cm}}
\toprule
项目 & 内容 \\
\midrule
需求描述 & 系统应允许 B 端用户通过 BearFRP 把本地 LLM-Wiki API 发布为 C 端可访问的公网 API URL。 \\
参与者 & B 端知识拥有者。 \\
前置条件 & B 端 LLM-Wiki API 健康检查通过；B 端可连接 S 端 frps；BearFRP 账号和代理创建能力可用。 \\
基本流程 & 用户在 B 端确认本地 LLM-Wiki API 地址和端口；登录或连接 BearFRP 控制面；选择 HTTP 子域名模式或 TCP 端口模式；系统创建代理并生成 frpc 配置或启动脚本；B 端启动 frpc 建立隧道；系统返回可分享的公网 LLM-Wiki API URL。 \\
异常流程 & 本地 API 不可达时禁止发布并提示先修复 LLM-Wiki；HTTP 通配域名未配置时提示切换 TCP 模式或补充域名；端口冲突、余额不足、配额不足或认证失败时，控制面返回明确原因；frpc 断开时，系统显示离线状态并允许用户重试。 \\
后置条件 & C 端可通过公网 URL 访问 B 端 LLM-Wiki API；B 端可关闭或回收该入口。 \\
验收标准 & 公网 URL 的健康检查和项目列表可访问；发布入口指向 LLM-Wiki API，不指向 B 端完整 OpenCode；关闭代理后公网入口不可继续访问。 \\
关联证据 & T-07 BearFRP published entry；T-06 auto-publish sidecar；桌面端联调流程。 \\
\bottomrule
\end{longtable}

\subsubsection{FR-04 S 端代理、隧道、端口/子域名和状态管理}
\begin{longtable}{>{\bfseries}p{2.4cm}p{11.8cm}}
\toprule
项目 & 内容 \\
\midrule
需求描述 & S 端应提供 BearFRP 控制面和 frps 转发能力，用于管理用户、代理、子域名/端口、frpc 配置、在线状态和基础资源约束。 \\
参与者 & S 端管理员、B 端知识拥有者。 \\
前置条件 & S 端公网服务器可运行 BearFRP backend、frps 和控制面数据库；端口和域名配置已准备。 \\
基本流程 & 管理员启动 BearFRP backend 和 frps；B 端用户创建 HTTP 或 TCP 代理；系统分配子域名或公网端口，生成 frpc 连接参数；frps 接收 frpc 连接并转发 C 端请求；控制面展示代理在线状态、连接时间、流量或连接数量；管理员可禁用异常代理或调整端口/子域名配置。 \\
异常流程 & 子域名已占用或端口冲突时，系统拒绝创建并提示更换；frps token 不匹配时，连接被拒绝并记录状态；代理离线或心跳超时时，控制面标记离线；资源超过限制时，系统应限流、拒绝新连接或提示升级配置。 \\
后置条件 & S 端只保存账号、代理和隧道状态，不保存知识正文和问答上下文。 \\
验收标准 & BearFRP API 可创建和查询代理；在线状态可观测；HTTP/TCP 模式至少一种可完成演示；管理员可停止异常代理。 \\
关联证据 & T-04 BearFRP user/proxy pytest；T-05 frps plugin and poller pytest；BearFRP 发布说明。 \\
\bottomrule
\end{longtable}

\subsubsection{FR-05 C 端添加远程知识库 API}
\begin{longtable}{>{\bfseries}p{2.4cm}p{11.8cm}}
\toprule
项目 & 内容 \\
\midrule
需求描述 & C 端用户应能在 WikiBridge Desktop 中添加 B 端分享的 LLM-Wiki API 地址，并完成健康检查、项目读取、Token 保存和项目选择。 \\
参与者 & C 端知识消费者。 \\
前置条件 & C 端可访问 B 端分享的公网 API URL；如 API 开启鉴权，用户持有 Token。 \\
基本流程 & 用户在“添加远程知识库”界面输入 API URL、名称和可选 Token；系统规范化 URL，并调用健康检查接口；系统调用项目列表接口；用户选择一个项目作为当前知识库；系统在 C 端本地保存远程知识库配置和最近检测状态。 \\
异常流程 & URL 格式错误时，界面提示修正；网络超时、DNS 失败或隧道离线时，系统提示远程不可达；401/403 时提示检查 Token；项目列表为空时，提示 B 端先构建或选择可分享项目。 \\
后置条件 & C 端远程知识库列表中出现该 API，并显示健康状态、当前项目和连接入口。 \\
验收标准 & C 端可保存、检测、切换和删除远程知识库；错误状态不会污染已有可用配置；Token 不以明文完整形式展示在日志和界面中。 \\
关联证据 & T-07 local entry OpenCode KB home；T-09 desktop project dashboard/compile/link screenshots。 \\
\bottomrule
\end{longtable}

\subsubsection{FR-06 C 端本地启动 OpenCode}
\begin{longtable}{>{\bfseries}p{2.4cm}p{11.8cm}}
\toprule
项目 & 内容 \\
\midrule
需求描述 & C 端应在本机启动 OpenCode，并以 KB 模式作为远程知识库对话入口。 \\
参与者 & C 端知识消费者。 \\
前置条件 & C 端已安装或打包 OpenCode sidecar；用户已保存模型供应商、模型名和 API Key；至少一个远程知识库可用。 \\
基本流程 & 用户保存或确认 C 端模型配置；点击连接远程知识库或启动 OpenCode；Desktop 启动本地 OpenCode 进程；OpenCode 以 KB 模式提供本地 Web 入口和健康检查；用户进入 KB Chat Session。 \\
异常流程 & OpenCode 二进制缺失或版本过旧时提示重新准备 sidecar；端口被占用时提示更换端口或停止旧进程；模型配置缺失时禁止进入正式问答并提示补充配置；OpenCode 启动后健康检查失败时，显示日志入口和重试按钮。 \\
后置条件 & C 端本地 OpenCode 可访问，页面标识 KB 模式，模型配置保存在 C 端。 \\
验收标准 & 本地 OpenCode 健康检查成功；页面含 KB 模式标记；C 端模型 API Key 不上传至 B 端或 S 端。 \\
关联证据 & T-03 KB mode meta and UI；T-07 local entry OpenCode KB home；OpenCode KB 模式说明。 \\
\bottomrule
\end{longtable}

\subsubsection{FR-07 注册并调用 llm-wiki MCP 工具}
\begin{longtable}{>{\bfseries}p{2.4cm}p{11.8cm}}
\toprule
项目 & 内容 \\
\midrule
需求描述 & C 端应把远程 LLM-Wiki API 注册为 OpenCode 可调用的 MCP 工具，并限制在 \code{llm_wiki_*} 知识库工具范围内。 \\
参与者 & C 端知识消费者、OpenCode、llm-wiki MCP server。 \\
前置条件 & C 端 OpenCode 已启动；远程知识库配置可用；MCP server 可运行；如远程 API 需要 Token，Token 已保存。 \\
基本流程 & Desktop 根据远程知识库生成 MCP server 配置；系统通过环境变量传入 API base URL 和 Token；OpenCode 注册 llm-wiki MCP；MCP server 暴露状态、项目、文件、读取、搜索和图谱等允许工具；OpenCode 在对话中调用 \code{llm_wiki_status}、\code{llm_wiki_projects}、\code{llm_wiki_read_file}、\code{llm_wiki_search} 等工具。 \\
异常流程 & MCP server 启动失败时提示检查 Node 版本、入口文件和依赖；远程 API 不可达时工具返回可读错误和 Desktop 修复提示；非 JSON 响应、API 错误和网络错误应被包装为可诊断工具错误；非白名单 MCP 或本地命令型 MCP 不得在 KB 模式下被任意添加。 \\
后置条件 & OpenCode 可以通过 MCP 访问远程知识库，但不能越过 LLM-Wiki API 直接扫描 B 端文件系统。 \\
验收标准 & MCP 测试通过；工具调用携带必要 Token；文件读取和搜索结果来自远程 LLM-Wiki API；非法 MCP 添加或危险工具调用被拒绝。 \\
关联证据 & T-02 LLM Wiki MCP tests；LLM-Wiki MCP server 说明；T-03 KB permission tests。 \\
\bottomrule
\end{longtable}

\subsubsection{FR-08 对话式读取、搜索和总结远程知识库}
\begin{longtable}{>{\bfseries}p{2.4cm}p{11.8cm}}
\toprule
项目 & 内容 \\
\midrule
需求描述 & C 端用户应能在本地 OpenCode 中用自然语言读取、搜索和总结 B 端分享的远程知识库。 \\
参与者 & C 端知识消费者。 \\
前置条件 & FR-05、FR-06、FR-07 已完成；C 端模型配置有效；远程知识库中存在可读内容。 \\
基本流程 & 用户在 KB Chat 中输入问题，例如要求读取某个 Markdown 文件或搜索关键词；OpenCode 根据问题选择 \code{llm_wiki_*} 工具；MCP server 请求远程 LLM-Wiki API；LLM-Wiki 返回文件内容、搜索结果或项目元数据；OpenCode 使用 C 端模型配置生成回答、摘要或引用说明。 \\
异常流程 & 未选择项目时，系统提示先选择远程知识库项目；文件不存在或路径不允许时，工具返回明确错误；内容过长时，工具可返回截断标记并提示分段读取；模型调用失败时，保留工具结果并提示用户检查 C 端模型配置。 \\
后置条件 & 用户获得基于远程知识库内容的回答；B 端只暴露受控 API 响应。 \\
验收标准 & C 端能完成至少一次文件读取、一次搜索和一次总结演示；回答不得要求 B 端执行终端命令或访问未分享路径。 \\
关联证据 & T-07 llm-wiki-file-content；T-07 llm-wiki-search-results；T-09 desktop local wiki reader。 \\
\bottomrule
\end{longtable}

\subsubsection{FR-09 KB 模式安全边界}
\begin{longtable}{>{\bfseries}p{2.4cm}p{11.8cm}}
\toprule
项目 & 内容 \\
\midrule
需求描述 & 系统必须在需求层明确 KB 模式安全边界：B 端不得暴露完整 OpenCode；C 端 OpenCode 只能通过受控 \code{llm_wiki_*} MCP 工具访问远程知识库；终端、本地命令、越权文件访问和非白名单 MCP 应被阻断。 \\
参与者 & B 端知识拥有者、C 端知识消费者、项目开发者、测试人员。 \\
前置条件 & C 端 OpenCode 以 KB 模式运行；远程知识库通过 LLM-Wiki API 接入；安全测试脚本可执行。 \\
基本流程 & 系统在 OpenCode 页面注入 KB 模式标识；文件访问限制在允许的知识库路径或远程 API 结果内；公开 Wiki 或远程知识库内容按只读方式消费；Shell、PTY、终端、配置修改、动态 MCP 添加和危险权限在 KB 模式下被拒绝或隐藏；日志隐藏 Token、API Key 和敏感路径。 \\
异常流程 & 用户直接调用受限 API 时返回 403 或等价拒绝结果；用户尝试路径穿越、绝对路径、软链接逃逸或其它用户目录访问时被拒绝；旧 sidecar 或旧脚本试图发布 B 端 OpenCode 时，应被视为不符合当前需求基线。 \\
后置条件 & C 端可以对话式消费知识库，但不能获得 B 端 Agent、本地命令或越权文件访问能力。 \\
验收标准 & KB 模式 meta 可检测；终端/PTY/配置/MCP 动态添加等高危接口被拒绝；B 端发布入口指向 LLM-Wiki API；安全测试和 XSS 检查通过。 \\
关联证据 & T-03 KB mode meta and UI；T-08 XSS no alert；OpenCode KB 模式说明中的已验证项。 \\
\bottomrule
\end{longtable}

\subsubsection{FR-10 部署、健康检查、联调和测试支撑}
\begin{longtable}{>{\bfseries}p{2.4cm}p{11.8cm}}
\toprule
项目 & 内容 \\
\midrule
需求描述 & 项目应提供可复现的部署、健康检查、联调和测试支撑，使需求验收不依赖口头说明。 \\
参与者 & 项目开发者、S 端管理员、测试人员。 \\
前置条件 & 代码仓库、Docker、Node、Rust/Cargo、测试脚本和演示数据按文档准备；部分测试需要 Node 20.19+。 \\
基本流程 & 使用 Docker Compose 或本地开发命令启动 LLM-Wiki、BearFRP、OpenCode 和必要网关；准备 \code{sample-wiki} 或等价演示数据；执行 LLM-Wiki API、BearFRP、OpenCode KB、MCP 和 Desktop 相关测试；记录命令输出、截图证据和剩余风险；将测试结果与 FR 编号关联。 \\
异常流程 & Node 版本不足时，脚本应提示所需运行时版本并终止对应自动化路径；Docker 或服务未启动时，黑盒检查应失败并说明不可达组件；真实 UI 截图采集失败时，应保留错误日志并重新执行截图脚本生成证据。 \\
后置条件 & 项目组能够用命令、截图和报告证明核心链路可运行。 \\
验收标准 & 测试脚本语法检查和空白检查通过；T-02/T-07 黑盒项目级检查通过；T-06 sidecar 检查通过；T-09 合同测试、Playwright 测试和截图采集通过。 \\
关联证据 & 测试命令与脚本输出整理；Playwright 人工测试截图说明；测试脚本说明。 \\
\bottomrule
\end{longtable}

\subsection{需求追踪矩阵}
\begin{longtable}{>{\centering\arraybackslash}p{1.6cm}p{3.6cm}p{4.8cm}p{4.2cm}}
\toprule
需求编号 & 关联用例 & 关联测试/证据 & 备注 \\
\midrule
FR-01 & 构建本地知识库 & \code{sample-wiki} fixture；T-02/T-07 & 演示数据规模以测试报告记录为准。 \\
FR-02 & API 健康、项目、文件、搜索 & T-02；T-07；文件/搜索截图 & API token 策略以实际配置为准。 \\
FR-03 & 发布 LLM-Wiki API & T-06；T-07 BearFRP 发布入口 & 不再发布完整 OpenCode。 \\
FR-04 & BearFRP 代理管理 & T-04；T-05；部署说明 & HTTP 子域名依赖通配域名。 \\
FR-05 & 添加远程知识库 & T-09 Desktop 截图；桌面端联调流程 & Desktop 截图编号见第 5 章。 \\
FR-06 & 启动本地 OpenCode & T-03；T-07 local entry & Node/sidecar 版本需满足打包要求。 \\
FR-07 & MCP 注册和工具调用 & T-02 MCP tests；T-03 & 工具命名采用 \code{llm_wiki_*}。 \\
FR-08 & 对话式消费 & T-07 文件/搜索截图；T-09 local wiki reader & 模型调用效果受 C 端模型配置影响。 \\
FR-09 & KB 安全边界 & T-03；T-08；KB 模式说明 & 旧脚本只作为历史资料，不作为当前需求。 \\
FR-10 & 部署与测试支撑 & 测试命令输出；Playwright 截图说明 & Node 22 补跑结果和截图证据已纳入测试报告。 \\
\bottomrule
\end{longtable}

\section{其它软件需求描述}

\subsection{性能要求}
\begin{longtable}{>{\bfseries}p{3.0cm}p{11.2cm}}
\toprule
编号 & 要求 \\
\midrule
NFR-PERF-01 & 在课程演示数据规模下，LLM-Wiki 健康检查、项目列表和文件列表应在可交互时间内返回；本地 P95 响应时间目标不超过 2 秒。 \\
NFR-PERF-02 & 通过 BearFRP/frp 访问远程 LLM-Wiki API 时，系统应尽量保持问答前工具调用可用；网络抖动时应返回超时提示，不应导致 Desktop 或 OpenCode 无响应。 \\
NFR-PERF-03 & C 端启动 OpenCode 和注册 MCP 的过程应给出进度或状态反馈；演示环境启动时间目标不超过 30 秒。 \\
NFR-PERF-04 & 搜索和文件读取结果过大时，API 或 MCP 工具应支持截断、分页或限制返回，避免单次工具调用阻塞对话。 \\
\bottomrule
\end{longtable}

\subsection{安全与隐私要求}
\begin{longtable}{>{\bfseries}p{3.0cm}p{11.2cm}}
\toprule
编号 & 要求 \\
\midrule
NFR-SEC-01 & B 端对外只发布 LLM-Wiki API，不发布完整 OpenCode 页面、终端、文件系统或 Agent 执行环境。 \\
NFR-SEC-02 & S 端只保存账号、代理、端口/子域名、隧道状态等控制面数据，不保存原始资料、Wiki 正文、向量索引和问答上下文。 \\
NFR-SEC-03 & C 端模型供应商、模型名和 API Key 只保存在 C 端本地，不上传给 B 端或 S 端。 \\
NFR-SEC-04 & Token、API Key、密码和 frps/frpc 鉴权信息不得以明文完整形式出现在日志、截图和普通界面中。 \\
NFR-SEC-05 & KB 模式必须拒绝终端执行、PTY 创建或连接、越权文件读取、配置修改和非白名单 MCP 动态添加。 \\
NFR-SEC-06 & 文件读取必须通过 LLM-Wiki API 或 KB 模式允许路径进行，路径穿越、绝对路径越权和软链接逃逸应被拒绝。 \\
NFR-SEC-07 & B 端发布入口应支持主动关闭、过期或管理员回收。滥用、异常流量或 Token 泄露时应能够停止对应代理。 \\
NFR-SEC-08 & 展示远程知识库内容时应防止脚本执行和弹窗注入；XSS 检查结果应纳入验收证据。 \\
\bottomrule
\end{longtable}

\subsection{可用性与可靠性要求}
\begin{longtable}{>{\bfseries}p{3.0cm}p{11.2cm}}
\toprule
编号 & 要求 \\
\midrule
NFR-REL-01 & LLM-Wiki API、BearFRP 代理、frp 隧道、远程知识库、OpenCode 和 MCP 注册状态应在界面或日志中可观测。 \\
NFR-REL-02 & 任一组件不可达时，系统应指明不可达组件和建议修复动作，例如检查 Token、端口、域名、隧道或 sidecar 版本。 \\
NFR-REL-03 & 发布入口、远程知识库配置和模型配置应在可恢复范围内本地持久化；失败重试不得删除已有可用配置。 \\
NFR-REL-04 & 黑盒检查和截图脚本应在环境不满足时给出明确版本或依赖提示，不应把环境限制误判为功能失败。 \\
\bottomrule
\end{longtable}

\subsection{兼容性与运行环境要求}
\begin{longtable}{>{\bfseries}p{3.0cm}p{11.2cm}}
\toprule
编号 & 要求 \\
\midrule
NFR-COMP-01 & WikiBridge Desktop 原型应面向主流桌面环境，课程验收阶段以 Linux 桌面环境和浏览器自动化截图为主要验证平台。 \\
NFR-COMP-02 & OpenCode、MCP server 和 Desktop 前端相关测试需要 Node 20.19+ 或 CI 中等价运行时；最终补跑环境采用 Node 22.13.1。 \\
NFR-COMP-03 & Docker Compose 演示环境应能启动 LLM-Wiki、BearFRP、OpenCode、nginx 和可选 frpc sidecar。 \\
NFR-COMP-04 & HTTP 子域名发布依赖公网域名和通配解析；未具备该条件时，系统应支持 TCP 端口模式完成演示。 \\
\bottomrule
\end{longtable}

\subsection{界面要求}
\begin{longtable}{>{\bfseries}p{3.2cm}p{11.0cm}}
\toprule
界面 & 需求 \\
\midrule
本地知识库/项目界面 & 展示本地项目、构建状态、文件或编译就绪状态；构建失败时给出可读错误。 \\
BearFRP 发布界面 & 展示本地 API 地址、代理类型、子域名/端口、隧道状态、公网 URL 和关闭入口。 \\
远程知识库界面 & 支持输入 API URL、名称、Token、健康检查、项目选择、连接、删除和错误提示。 \\
模型配置界面 & 支持 C 端填写模型供应商、模型名、API Key，并说明该配置保存在 C 端。 \\
OpenCode 入口界面 & 展示本地 OpenCode 启动状态、KB 模式状态、进入对话入口和运行日志。 \\
MCP 连接界面 & 展示 MCP 注册状态、允许工具列表和失败原因。 \\
错误与日志界面 & 错误信息应面向用户可理解，不只暴露底层堆栈；敏感信息应脱敏。 \\
\bottomrule
\end{longtable}

\subsection{进度要求}
\begin{longtable}{>{\bfseries}p{3.2cm}p{11.0cm}}
\toprule
项目 & 要求 \\
\midrule
课程阶段主链路 & 优先完成 B 端 LLM-Wiki API 发布、C 端远程知识库添加、OpenCode 启动、MCP 注册和对话式读取/搜索。 \\
文档交付 & 需求规格说明书、软件设计说明书、测试报告、计划任务书和个人总结材料应保持术语一致。 \\
最终交付基线 & 2026年6月26日形成 V1.0 需求、SDS V1.2 设计和 V1.0 测试材料；T-09 Desktop 截图与 Node 22 补跑结果纳入最终验收证据。 \\
\bottomrule
\end{longtable}

\subsection{交付要求}
\begin{longtable}{>{\bfseries}p{3.2cm}p{11.0cm}}
\toprule
类型 & 交付物 \\
\midrule
程序 & WikiBridge Desktop、LLM-Wiki API、BearFRP/frps 集成、OpenCode KB 模式、llm-wiki MCP server、Docker Compose 和测试脚本。 \\
文档 & 本需求规格说明书、计划任务书、软件设计说明书、测试文档/测试报告、部署说明、Desktop 交接说明和演示说明。 \\
证据 & 测试命令输出、Playwright 截图、API 证据页、架构演进说明和最终答辩截图。 \\
\bottomrule
\end{longtable}

\subsection{验收要求}
\begin{enumerate}
  \item 文档验收：本文档应覆盖 GB/T 8567-88 和课程参考 SRS 的核心章节，包括引言、一般性描述、功能需求、其它需求和软件原型。
  \item 功能验收：FR-01 至 FR-10 每项均应具备参与者、前置条件、基本流程、异常流程和验收标准。
  \item 链路验收：B 端构建 Wiki、发布 LLM-Wiki API，C 端添加远程知识库、本地启动 OpenCode、注册 MCP，并完成读取/搜索/总结演示。
  \item 安全验收：不得通过公网暴露 B 端完整 OpenCode；C 端 KB 模式应拒绝终端、本地命令、越权文件和非白名单 MCP。
  \item 测试验收：以测试命令与脚本输出整理、测试分析报告和 CI 结果为依据，记录通过、失败和环境差异。
  \item 截图验收：第 5 章列出的截图证据已按测试项编号，可用于最终提交和答辩演示。
\end{enumerate}

\section{软件原型与演示证据}

\subsection{原型界面文字描述}
当前 WikiBridge 原型以桌面端工作台和本地 OpenCode KB 页面为主要入口：
\begin{itemize}
  \item Desktop 本地知识库界面：展示项目列表、构建状态、编译就绪状态和本地 Wiki 阅读入口。该界面用于证明 B 端可以准备和检查本地知识库。
  \item Desktop BearFRP/访问连接界面：展示发布 API、代理状态、公网 URL 和联调错误。该界面用于证明 B 端分享的是 LLM-Wiki API。
  \item Desktop 远程知识库界面：允许 C 端粘贴 B 端分享地址、填写 Token、执行健康检查、读取项目列表并连接远程知识库。
  \item Desktop OpenCode 面板：负责启动本地 OpenCode，显示运行状态，并进入 KB Chat。
  \item OpenCode KB 页面：页面带有 KB 模式标识，隐藏或拒绝终端、危险文件访问和动态 MCP 等能力，作为 C 端对话式消费入口。
  \item API 证据页：当真实 UI 截图受环境限制时，可使用同源 API 返回数据生成证据页，证明文件内容读取、搜索结果和安全检查通过。
\end{itemize}

\subsection{已完成截图证据}
\begin{longtable}{>{\centering\arraybackslash}p{1.3cm}p{6.2cm}>{\centering\arraybackslash}p{1.6cm}>{\centering\arraybackslash}p{2.0cm}p{2.7cm}}
\toprule
测试项 & 截图文件 & 当前结果 & 证据来源 & 对应需求 \\
\midrule
T-07 & 本地入口 OpenCode KB 页面截图 & PASS & real-ui & FR-06、FR-09 \\
T-03 & KB 模式页面与 meta 标识截图 & PASS & real-ui & FR-06、FR-09 \\
T-07 & LLM-Wiki 知识库页面截图 & PASS & api-evidence & FR-01、FR-02 \\
T-07 & LLM-Wiki 文件内容页面截图 & PASS & api-evidence & FR-02、FR-08 \\
T-07 & LLM-Wiki 搜索结果页面截图 & PASS & api-evidence & FR-02、FR-08 \\
T-08 & XSS 样例页面截图 & PASS & api-evidence & FR-09 \\
T-07 & BearFRP 发布入口页面截图 & PASS & real-ui & FR-03、FR-04 \\
T-09 & 桌面端知识库项目仪表盘截图 & PASS & real-ui & FR-01、FR-05 \\
T-09 & 桌面端编译准备状态截图 & PASS & real-ui & FR-01 \\
T-09 & 桌面端访问连接状态截图 & PASS & real-ui & FR-03、FR-05 \\
T-09 & 桌面端消费端远程知识库入口截图 & PASS & real-ui & FR-01、FR-08 \\
\bottomrule
\end{longtable}

\subsection{演示脚本}
\begin{enumerate}
  \item B 端启动 LLM-Wiki，并确认 \code{/api/v1/health} 和项目列表可访问。
  \item B 端通过 BearFRP 创建代理，获得公网 LLM-Wiki API URL。
  \item C 端打开 WikiBridge Desktop，保存自己的模型配置。
  \item C 端在远程知识库界面粘贴公网 API URL，执行健康检查并选择项目。
  \item C 端连接该远程知识库，Desktop 启动本地 OpenCode 并注册 llm-wiki MCP。
  \item C 端在 KB Chat 中提问，例如“读取知识库里的概览文件并总结主要内容”或“搜索 WikiBridge 相关内容”。
  \item 演示安全边界：确认页面处于 KB 模式，终端/PTY/动态 MCP 等高危能力不可用，B 端没有向公网暴露完整 OpenCode。
\end{enumerate}

\end{document}
